% Verification-First Legal RAG: A Defense-in-Depth Audit Pipeline and a
% Multilingual Swiss-Law Benchmark
%
% arXiv submission v2 — paper.tex skeleton, populated in Step 5.
%
% Style: NeurIPS-like single-column LaTeX (we'll switch to ACL or
% specific arXiv-friendly class once we confirm the target during Step 1).
% This skeleton is deliberately minimal so the structure is the
% review surface, not the formatting.

\documentclass[11pt,a4paper]{article}

\usepackage[utf8]{inputenc}
\usepackage[T1]{fontenc}
\usepackage[margin=2.5cm]{geometry}
\usepackage{microtype}
\usepackage{booktabs}
\usepackage{graphicx}
\usepackage{hyperref}
\usepackage{xcolor}
\usepackage{authblk}
\usepackage{abstract}

\hypersetup{
    colorlinks=true,
    linkcolor=blue!50!black,
    urlcolor=blue!50!black,
    citecolor=blue!50!black,
}

\title{Verification-First Legal RAG:\\
A Defense-in-Depth Audit Pipeline and\\
a Multilingual Swiss-Law Benchmark}

\author[1]{Jonas Hertner}
\affil[1]{OpenCaseLaw, Zurich, Switzerland \\ \texttt{team@jonashertner.com}}

\date{\today}

\begin{document}

\maketitle

\begin{abstract}
% [Abstract draft A — locked in Step 1, populated in Step 5 once
% per-rail ablation numbers are in.]
%
% TODO: paste from 01_title_and_abstract.md once Step 1 questions
% are answered and Draft A vs B vs C is chosen.
\end{abstract}

% ──────────────────────────────────────────────────────────────────────

\section{Introduction}
\label{sec:intro}

% [§1 — open with Dahl 2024 + Magesh 2024 numbers; introduce the
% rail decomposition as the conceptual move; state the two
% contributions; close with paper roadmap.]

% Length budget: 0.75 page.

% Outline:
% - Sentence 1: Magesh / Dahl finding (58–82 % / 17–33 % hallucination).
% - Sentence 2: Why this matters specifically for legal practice
%   (Sorgfaltspflicht / professional liability).
% - Sentence 3: The conceptual move — decompose verification by
%   error class.
% - Paragraph 2: The two contributions (C1, C2 from
%   02_contribution_table.md).
% - Paragraph 3: Roadmap (§3 architecture, §4 benchmark, §5 results,
%   §6 discussion, §7 conclusion).

% ──────────────────────────────────────────────────────────────────────

\section{Related Work}
\label{sec:related}

% [§2 — Tier 1 / Tier 2 / Tier 3 / Tier 4 from 03_related_work_map.md
% prose-ified into 4 paragraphs.]

% Length budget: 0.75 page.

% Closing paragraph: the "positioning paragraph" from
% 03_related_work_map.md.

% ──────────────────────────────────────────────────────────────────────

\section{The Five-Rail Audit Pipeline}
\label{sec:rails}

% [§3 — the methodological contribution.]

% Length budget: 1.5 pages.

% Outline:
% - §3.1 Architecture overview (figure 1: pipeline diagram)
% - §3.2 Rail 1: case-citation existence (regex parse, exact-match
%   lookup, pinpoint validation against structured Erwägungen +
%   full-text fallback)
% - §3.3 Rail 2: statute-reference resolution (regex parse, statute-DB
%   lookup, abbreviation guard)
% - §3.4 Rail 3: verbatim-quote source-matching (regex parse,
%   whitespace + punctuation normalisation, source-pool from cited
%   decisions + cited statutes)
% - §3.5 Rail 4: decision-date sanity (date adjacency parse,
%   match against stored decision_date)
% - §3.6 Rail 5: opt-in proposition grounding (preceding-claim
%   extraction, batched Sonnet judge, JSON verdict array)
% - §3.7 Cost and latency budget (Table 1: per-rail cost; Figure 2:
%   latency histogram on 100 sample drafts)

% ──────────────────────────────────────────────────────────────────────

\section{Swiss Legal RAG Bench v0.2}
\label{sec:benchmark}

% [§4 — the resource contribution.]

% Length budget: 1.5 pages.

% Outline:
% - §4.1 Design rationale (extension of Butler & Butler 2026 with
%   cross-lingual leaf)
% - §4.2 Question construction (expert-annotated, balanced across
%   DE/FR/IT × major legal areas; 30 questions)
% - §4.3 Evaluation methodology (5-leaf decomposition: correctness ×
%   groundedness × retrieval accuracy × cross-language)
% - §4.4 Statistics (Table 2: distribution by language / area /
%   difficulty; Table 3: claim-type distribution)

% ──────────────────────────────────────────────────────────────────────

\section{Experiments}
\label{sec:experiments}

% [§5 — the empirical core.]

% Length budget: 2 pages.

% Outline:
% - §5.1 Setup (live OpenCaseLaw stack as retriever; Claude Sonnet 4.6
%   as generator and judge; Step 4 produces the actual numbers)
% - §5.2 Per-rail ablation (Table 4: 6 audit configurations × 30
%   questions; per-rail catch rate per error class)
% - §5.3 Comparison against Butler & Butler 2026 baseline (Table 5:
%   our Swiss numbers vs their Australian numbers, with same metrics)
% - §5.4 Cross-lingual subset analysis (where retrieval succeeds in
%   the question's language vs cross-language; Figure 3: language
%   confusion matrix)
% - §5.5 Cost / latency tradeoff (4 deterministic rails ~3 ms;
%   judge rail ~2 s, ~$0.005)

% ──────────────────────────────────────────────────────────────────────

\section{Discussion}
\label{sec:discussion}

% [§6 — limitations and generalisation.]

% Length budget: 0.75 page.

% Outline:
% - §6.1 When the deterministic rails miss (paraphrased quotes;
%   indirect citations; multi-source propositions)
% - §6.2 When the judge rail miscalls (long claims with multiple
%   sub-propositions; claims requiring inference beyond the
%   retrieved text)
% - §6.3 Generalisation to other jurisdictions (each rail's
%   prerequisite; what a non-Swiss adaptation would change)
% - §6.4 Limitations (30-question benchmark size; live-corpus
%   dependency in some experiments; reliance on Claude as both
%   generator and judge for cost reasons)

% ──────────────────────────────────────────────────────────────────────

\section{Conclusion}
\label{sec:conclusion}

% [§7 — concise.]

% Length budget: 0.25 page.

% Restate the two contributions; point to code (MIT) + benchmark
% (CC0) + corpus (CC0); thank the OnlineKommentar.ch and
% entscheidsuche.ch communities; gesture at v0.3 of the benchmark
% (100 questions) and the open question of automated benchmark
% growth.

% ──────────────────────────────────────────────────────────────────────

\section*{Data Availability}

% Audit pipeline: \url{https://github.com/jonashertner/caselaw-repo-1}
% Benchmark: \url{https://huggingface.co/datasets/voilaj/swiss-legal-rag-bench}
% Corpus: \url{https://huggingface.co/datasets/voilaj/swiss-caselaw}
% Live MCP endpoint: \url{https://mcp.opencaselaw.ch}
% Frozen paper-release snapshot (DOI): % [Zenodo DOI from Step 2]

% ──────────────────────────────────────────────────────────────────────

\bibliographystyle{plain}
\bibliography{bib/references}

\end{document}
